
\chapter{Le département des Arts du spectacle de la BnF}

Le département des Arts du Spectacle est l'un des sept départements spécialisés de la Bibliothèque nationale de France. Celui-ci a pour particularité de rassembler des documents en fonction de leur sujet, et non de leur nature. Le département conserve donc des documents écrits, qu'ils soient imprimés ou manuscrits, des documents iconographiques, et même des documents sonores, audiovisuels ou des objets qui ont tous pour sujet les arts du spectacle\footcite[p. 8-9]{giteau_1995}.

\section{La création du département}

Le département est officiellement créé le 1er janvier 1976. Il est notamment constitué à partir de la collection d'Auguste Rondel, et sa création fait suite à un demi-siècle d'enrichissements. En effet, en 1935, juste après la mort de Rondel, la bibliothèque de l'Arsenal, qui conserve la collection, est rattachée à la Bibliothèque nationale de France. Cette décision s'explique en partie par la volonté de Julien Cain, alors administrateur général de la Bibliothèque nationale et introduit dans les milieux du théâtre, et témoigne d'un souci de reconnaître un patrimoine artistique et documentaire relatif aux arts du spectacle. En 1953, il confie l'avenir de la collection à André Veinstein, qui a pour mission de la réactualiser. Ce dernier se lance donc dans une collecte de documents, afin d'accroître cette collection. Dans cette initiative, il est guidé par la volonté constante d'associer les instances professionnelles du spectacle et les acteurs de ce milieu 
à une réflexion sur la mémoire du théâtre, en lien avec la création théâtrale. 

Cette préfiguration du Département des Arts du spectacle se concrétise en 1958, lorsque l'entrée de la collection Craig signe le début d'une section \enquote{théâtre} dans les locaux de la Bibliothèque nationale, rue de Richelieu, rattachée aux services de l'Administrateur général et gérée par un conservateur. Cela se poursuit en 1965, lorsque des mesures administratives internes décident entre autres de l'attribution élargie et nominale aux collections théâtrales du dépôt légal pour les ouvrages et périodiques français\footcite[p. 250]{giteau_2005}.

La création du département, qui officialise le regroupement des collections de la Bibliothèque nationale de France à partir de la collection Rondel, est rapidement suivie par celle de la maison Jean Vilar à Avignon en 1979, qui marque l'ouverture du département au spectacle vivant. Les accroissements se poursuivent ensuite dans le département, dans la continuité de la collection Rondel. Les entrées sont partagées 
entre documents apportés par des comédiens et des artistes ou par leurs proches, et fonds importants d'auteurs dramatiques, comme le fonds volumineux de Sacha Guitry, acquis en 1995. Le département conserve également des documents liés aux décorateurs et aux créateurs de costumes ayant travaillé en France de la fin du XIX\textsuperscript{e} siècle à aujourd'hui. Les autres arts du spectacle, comme le cirque ou la danse, sont aussi représentés dans le département, et ont régulièrement été mis en valeur par des expositions dans la galerie Colbert, mise à disposition du Département entre 1985 et 1997. Sont également entrés au département des fonds concernant des grands réalisateurs de cinéma, comme Jean Grémillon en 1977, René Clair en 1984 ou encore Abel Gance en 1993 et 1995\footcite[p. 253-257]{giteau_2005}. 


\section{Le département aujourd'hui}

Dirigé par Joël Hutwohl depuis 2008, le département des Arts du spectacle est réparti en quatre services. Trois d'entre eux, à savoir le Service de l’iconographie et de la documentation, le Service des archives et des imprimés et le Service de la conservation et de la communication, se situent à Paris, et la Bibliothèque de la Maison Jean Vilar, qui constitue le quatrième est, nous l'avons vu précédemment, située à Avignon. Ses collections sont décrites dans deux catalogues, \enquote{BnF - Catalogue général} et \enquote{BnF - Archives et manuscrits}, et le département contribue régulièrement à des expositions sur les différents sites de la BnF\footcite{sitebnf}, en plus de prêter pour des expositions externes à l'établissement. 

En somme, la collection Rondel est centrale dans l'histoire du département des Arts du spectacle, d'où un besoin d'en connaître mieux l'histoire et les provenances, à l'heure où cette question est au coeur des préoccupations de nombreuses institutions patrimoniales. 






